% Part: computability
% Chapter: computability-theory
% Section: reducibility

\documentclass[../../../include/open-logic-section]{subfiles}

\begin{document}

\olfileid{cmp}{thy}{red}
\olsection{न्यूनीकरणक्षमता}

\begin{explain}
आता आपल्याला माहीत आहे की $K_0$ हा किमान एक संच संगणनक्षमपणे
प्रगणनीय आहे, पण संगणनक्षम नाही. असे इतरही संच असतील हे उघड
आहे. प्रत्येक वेळी मूलभूत व्याख्यांपासून पुन्हा सुरुवात न करता
इतर संचांचे असे गुणधर्म सिद्ध करण्याची प्रभावी पद्धत
न्यूनीकरणक्षमता देते.

सर्वसाधारणपणे, $A$ या संचाचे~$B$ कडे ``न्यूनीकरण'' म्हणजे
!!{element}s हे~$B$ मध्ये आहेत की नाहीत या प्रश्नांची उत्तरे
रूपांतरित करून !!{element}s हे~$A$ मध्ये आहेत की नाहीत या प्रश्नांची उत्तरे
मिळवण्याची पद्धत. येथे आपण ``अनेक-एक न्यूनीकरणक्षमता''
विचारात घेऊ; वेगवेगळे गुणधर्म असणाऱ्या न्यूनीकरणक्षमतेच्या इतरही
अनेक संकल्पना आहेत. संगणकीय गुंतागुंतीच्या अभ्यासातही अशा
संकल्पना महत्त्वाच्या असतात; तेथे कार्यक्षमतेचाही विचार करावा
लागतो. उदाहरणार्थ, एखादा संच NP मध्ये असेल आणि प्रत्येक
NP समस्या त्याच्याकडे न्यूनीत करता येत असेल, तर त्याला
``NP-पूर्ण'' म्हणतात. येथे पुढे येणाऱ्या न्यूनीकरणासारखेच
न्यूनीकरण तेथे वापरतात, पण त्यासाठी ते बहुपदी वेळेत संगणित
करता येण्याची अतिरिक्त अट असते.

न्यूनीकरणाची कल्पना आपण आधीच अप्रत्यक्षपणे वापरली आहे.
$K$ हा संच पुढीलप्रमाणे ठरवा:
\[
K = \Setabs{x}{\cfind{x}(x) \fdefined},
\]
म्हणजे $K = \Setabs{x}{x \in W_x}$. थांबण्याची समस्या
अनिर्णेय असल्याचा आपला पुरावा (\olref[hlt]{thm:halting-problem})
सर्वांत थेटपणे $K$ संगणनक्षम नाही हे दाखवतो. $K_0$ हा
पुढील संच आहे, हे आठवा:
\[
K_0 = \Setabs{\tuple{e, x}}{\cfind{e}(x) \fdefined },
\]
म्हणजे $K_0 = \Setabs{\tuple{e,x}}{x \in W_e}$.
%READERNOTE{OLCMP-021 — गोठवलेल्या इंग्रजीत K_0 ची पहिली व्याख्या जोडी (e,x) वापरते, पण लगेचची समतुल्य व्याख्या (x,e) अशी उलट क्रमाने लिहिते. W_e मध्ये x असल्याचा अर्थ पहिल्या व्याख्येतील (e,x) जोडीशी जुळतो; म्हणून येथे तीच क्रमरचना ठेवली आहे.}
$K$ च्या असंगणनक्षमतेच्या कोणत्याही पुराव्याचा विस्तार करून
$K_0$ ची असंगणनक्षमता दाखवणे सोपे आहे: $K_0$ संगणनक्षम असता,
तर !!a{element}~$x$ हा $K$ मध्ये आहे की नाही हे फक्त
$\tuple{x, x}$ ही जोडी $K_0$ मध्ये आहे का ते विचारून
निर्णीत करता आले असते. $x$ ला $\tuple{x, x}$ कडे नेणारे
फलन~$f$ हे $K$ चे $K_0$ कडे \emph{न्यूनीकरण} करण्याचे
उदाहरण आहे.
\end{explain}

\begin{defn}
$A$ आणि $B$ हे नैसर्गिक संख्यांचे संच असू द्या.
संगणनक्षम फलन~$f\colon \Nat \to \Nat$ हे $A$ चे~$B$
कडे \emph{अनेक-एक न्यूनीकरण} असते जर आणि फक्त जर
प्रत्येक नैसर्गिक संख्या~$x$ साठी
\[
x \in A \quad \text{तेव्हाच} \quad f(x) \in B.
\]
असे न्यूनीकरण~$f$ अस्तित्वात असल्यास, $A$ हे~$B$ कडे
\emph{अनेक-एक न्यूनीत करता येते} असे म्हणतो; हे
$A \leq_m B$ असे लिहितो. $A$ चे~$B$ कडे आणि उलटही
अनेक-एक न्यूनीकरण शक्य असल्यास, $A$ आणि $B$ हे
\emph{अनेक-एक सममूल्य} आहेत असे म्हणतो; हे
$A \equiv_m B$ असे लिहितो.
\end{defn}

वरील व्याख्येतील फलन $f$ एकास-एक असेल, तर $A$ हे~$B$
कडे \emph{एकास-एक न्यूनीत करता येते} असे म्हणतात.
पुढे दिलेली बहुतेक न्यूनीकरणे ही अधिक मजबूत अट पूर्ण करतात;
पण आपण ती बाब वापरणार नाही.

\begin{digress}
एकास-एक न्यूनीकरणक्षमता ही अनेक-एक न्यूनीकरणक्षमतेपेक्षा खरोखर
अधिक मजबूत अट आहे; हे खरे आहे, पण अजिबात उघड नाही. म्हणजे,
असे अनंत संच $A$ आणि~$B$ आहेत की $A$ हे~$B$ कडे अनेक-एक
न्यूनीत करता येते, पण~$B$ कडे एकास-एक न्यूनीत करता येत नाही.
\end{digress}

\end{document}
